% !TeX root = ../main.tex

% Plan général de l'introduction : 

% \begin{itemize}
% 	\item Présentation de la gestion des collections : qu'est ce que c'est ? à quoi ça sert ? Pourquoi le faire ? Brève histoire \\ FAIT 2 pages
	
% 	\item Limites actuelle de la gestions des collections et pourquoi l'IA peut être une solution à ces limites, à quel but ?\\ FAIT 2,5 pages
	
% 	\item Présentation de la mission de stage répondant à l'idée d'intégrer l'IA dans la gestion des collections\\ FAIT 1 page
% 	\subitem Présentation de l'entreprise de stage 
% 	\subitem Logiciel Skin en bref
% 	\subitem Missions de stages \\
	
% 	\item Comment je vais tenter d'apporter des pistes et réponses à ce problème. 
% 	\subitem Problématique 
% 	\subitem Annonce de \terme{plan}
% \end{itemize}


		La conservation du patrimoine et la gestion des collections patrimoniales soulèvent des enjeux politiques, sociaux, culturels et identitaires. Avec l'arrivée du numérique, ces notions n'ont cessé d'évoluer comme le souligne Marcello La Guardia et Mila Koeva : 
		
		\begin{displayquote}
			\guillemetleft \textit{Recent years have been characterised by a profound transformation in the field of Cultural Heritage, shaped by interdisciplinary methodologies, digital innovation, and the demand for more inclusive, sustainable management approaches.} \guillemetright \footcite{laguardiaDigitalInnovationDocumentation2025}
		\end{displayquote}
		 
		
		% L'essor des technologies numériques et méthodologies multidisciplinaires basées, pour tout ou partie, sur l'utilisation de l'\gls{IA} amène à repenser la gestion de nos collections physiques et numériques. Notamment face à de nouveaux enjeux issus de la numérisation et de la diffusion des collections à travers un nouveau moyen de communication : internet. De telles innovations amènent de profondes restructurations de ces notions.
		
		Dès lors nous devons définir ce que nous entendons par \terme{patrimoine} et \terme{gestion des collections}. Quels sont les objets, lieux et pratiques encadrés par ces notions ? Quels en sont les problématiques et limites particulières ? Et enfin, comment patrimoine, gestion des collections et numériques s'imbriquent ?
		\newline
		
		Patrimoine et gestion des collections sont des notions variant en fonction des contextes institutionnels et culturels. En France, elles prennent racines au XVIII\textsuperscript{ème} siècle lors de la Révolution française au travers des confiscations révolutionnaires. 
		
		Entre 1790 et 1793, les biens monarchiques et ecclésiastiques sont placés sous patronage de la Nation. Certains sont revendus pour pallier la crise financière en vigueur, d'autres sont précieusement conservé.\footcite[p.21]{varryConfiscationsRevolutionnaires2009} Mais la position révolutionnaire n'est pas unanime : si certains cherchent à conserver pour la Nation, d'autres souhaitent détruire ces symboles de l'Ancien Régime ce qui entraîne des destructions, dégradations et vols. \footcite[p.134]{mellotConfiscationsRevolutionnairesHistoire2019}.
		
		Le patrimoine s'affirme en réaction à ces actes au travers du rapport du 31 août 1794 sur les exactions révolutionnaires de l'Abbé Henri Grégoire \textit{(04 décembre 1750 - 28 mai 1831)}. Il y dénonce le \gls{vand}\footcite{sprigathVandalismeRevolutionnaire179217941980a} des \guillemetleft [...] mobiliers appartenants à la Nation [...] \guillemetright \footcite{LabbeGregoire31}. L'historien Dominique Poulot approfondit : en réalité on souhaite préserver ce qui est justifié par \guillemetleft \space l'appel à l'avenir \guillemetright pour un idéal associant \guillemetleft légitimité nouvelle, utilité savante, et conservation raisonnée. \guillemetright \footcite[{p.183-184}]{LabbeGregoire31}.
		
		Le projet de \textit{Bibliographie Universelle de la France} \footcite{mellotConfiscationsRevolutionnairesHistoire2019} de 1790, première grande entreprise de catalogage national, ainsi que la recherche constante de lieu de conservation pour ces biens, témoignent de cette naissance de la notion de Patrimoine et de ses particularités françaises. 
		 
		L'archiviste paléographe Olivia Brissaud soulève dans son article de 2013\footcite{brissaudLelaborationPrincipeDinalienabilite2013/dec./10} les particularités patrimoniales françaises suivantes : d'abord le patrimoine est \guillemetleft \gls{ina} \guillemetright  et \guillemetleft \gls{inv} \guillemetright. Ensuite, il relève d'une responsabilité commune. De plus, il devient le fondement identitaire et culturel de la Nation. Enfin, l'État s'octroie la souveraineté juridique sur le patrimoine à l'aide d'un appareil bureaucratique dédié, placé sous sa juridiction directe, devenu depuis le Service des Musées de France.
		
		Ainsi, le patrimoine est une notion culturelle, aux racines profondément politiques et juridiques, constituant le socle identitaire national. il se matérialise dans des collections d'objets, dont la constitution a permis l'émergence de la notion de patrimoine.  \footcite[p.17]{ministeredelacultureGuideGestionDocuments2021}. Leur gestion relève directement de la juridiction de l'État français chargé de répondre aux enjeux politiques, juridiques et culturels associés.\newline
		
		Cependant, le patrimoine est un objet vivant qui ne cesse de se diversifier notamment avec l'inclusion de l'\gls{imgani}\footnote{L'invention de \guillemetleft l'image animée \guillemetright \space en tant que tel est difficile à dater. Le brevetage en 1895 du cinématographe des frères Auguste et Louis Lumière consacre cette nouvelle typologie documentaire.} en 1936\footnote{Création de la Cinémathèque française} puis celle de la photographie\footnote{Bien que Louis Daguerre grâce à son Daguerréotype, le premier \guillemetleft appareil photo \guillemetright, a été considéré comme l'inventeur de la photographie, son invention remonterait aux travaux de Nicéphore Niépce et de son héliographe. Voir \cite{willfriedbaatzPhotography1997}} en 1964\footnote{Création du Musée Français de la photographie}. D'autres typologies d'objets entrent au patrimoine comme les espaces naturels et le patrimoine immatériel en 2003, avec la \textit{Convention pour la sauvegarde du patrimoine culturel} promulguée par l'Organisation des Nations unies pour l'éducation, la science et la culture appelé l'\gls{UNESCO}
		
		Pour la gestion des collections, depuis la tentative de normalisation en 1890\footcite{denoelLapportLeopoldDelisle2019}, les pratiques de catalogage n'évoluent pas beaucoup malgré la diversification des collections. C'est l'informatique qui apporte un nouveau souffle à ce domaine.\newline 
		
		En 1970, la première base de données informatique patrimoniale est créée afin de constituer un premier catalogue des collections mutualisées : c'est JOCONDE qui est destinée aux Musées de France\footcite[p. 2]{verdierSystemeDescriptifObjets1999}. En parallèle, des systèmes de gestion des collections, ou \terme{SGC} pour Système de Gestion des Collections, sont développés : c'est le cas du projet GRIPHOS \textit{Museum Computer Network} à New York en 1967 qui permet de créer un inventaire et suivi informatique des collections physiques dans des musées. 
		
		La réduction de la taille des ordinateurs permet d'équiper informatiquement les institutions en masse. Ainsi en 1985 est installé le \terme{SGC} \guillemetleft Micromusées \guillemetright \space à l'échelle nationale sous l'impulsion du Service des Musées de France.\footcite[p.2]{dejantiMuseologieInformatiqueLoeuvre2018}. L'invention d'internet en 1983 et du \textit{World Wide Web} en 1989 permettent de faire de JOCONDE et Micromusée des outils nationaux pour une mutualisation complète des collections et pratiques.
		
		Des organisations comme l'\gls{IFLA} ou le \gls{CIDOC}, profitent de l'essor de l'informatique pour normaliser mondialement les pratiques de gestion des collections. Ainsi, des normes de description, leurs traductions informatiques appelées \guillemetleft standards d'encodage \guillemetright \space et des modèles conceptuels de gestion des collections sont créés.
		
		Le patrimoine devient un bien commun à normaliser dans sa description afin d'en faciliter l'échange et la diffusion nationale et mondiale. Pour ces deux raisons, l'État français incite à informatiser les musées dans un premier temps, puis toute institution possédant des collections patrimoniales.\newline
		
		La dématérialisation des catalogues amène à penser une solution de conservation numérique des collections pensée pour concilier conservation et communication. Alors apparaît la numérisation, processus qui aboutit à la création d'un \guillemetleft jumeau numérique \guillemetright \footcite{bergeotJumeauNumeriqueDans2025} afin de protéger l'objet physique tout en offrant la possibilité de le consulter\footnote{Pour plus d'informations : \url{https://www.culture.gouv.fr/thematiques/innovation-numerique/soutenir-la-creation-numerique-et-l-innovation/numerisation-des-contenus-culturels}}. Pour certaines typologies telles que l'\gls{imgani}, le numérique devient le moyen de création et de conservation privilégié. Cela s'explique par les avancées technologiques dans ce domaine ainsi que la popularisation de la vidéo numérique.  %comme le montre les aides mises en places par le Ministère de la Culture pour la numérisation des collections dès 1996 \footcite{bequetNumerisationPatrimoineDocumentaire2000}.
		
		La numérisation du patrimoine amène en 1997 la \gls{BNF} à ouvrir son site de consultation de documents numérisés : Gallica. C'est le premier service web dédié à la consultation de documents patrimoniaux numérisés. Puis en 2009, l'UNESCO reconnaît aux objets numériques, natifs ou numérisés, une valeur patrimoniale \footcite{CharteConservationPatrimoine}. Dès lors, le patrimoine numérique connaît une croissance sans précédent : par exemple la \gls{BNF} possédait environ trois-cent cinquante mille documents numériques en 2007 contre plus de dix millions en 2026.\newline
		
		Les diverses évolutions du patrimoine et de la gestion des collections que nous avons retracées imposent de nouvelles limites et difficultés. Les collections patrimoniales sont aujourd'hui doublement hétérogènes : plusieurs typologies documentaires y coexistent et les objets physiques côtoient leurs équivalents numériques \footnote{Voir la notion de \guillemetleft jumeau numérique \guillemetright dans  \cite{bergeotJumeauNumeriqueDans2025}}. Une seule collection appelle une gestion plurielle à travers l'application de différentes normes, standards et modèles conceptuels. À cela s'ajoute l'intensification de l'expansion des collections patrimoniales notamment due au numérique. 
		
		La production intensive de données numériques  participe à l'explosion quantitative des données que l'on apppelle la \gls{bgdt}. À titre d'exemple, en 2023 l'INA stocke plus de 60 pétaoctets de données, soit 60 000 000 de gigaoctets de données. \footcite{joannaStockerDonneesPiste2023}. Avec de telles quantités de données, deux enjeux se renouvèlent : celui du stockage, qui engage des questions d'infrastructures, et celui de l'accès, qui amène à repenser les moyens d'accès aux collections.
		
		Si l'informatique a apporté des solutions de gestion des collections, la diversité des modalités et pratiques ainsi que l'importance desdites collections appellent d'autres solutions. Depuis 2020, le domaine de l'\gls{IA} semble proposer des solutions face aux défis posés par les collections patrimoniales aujourd'hui. 
		
		On entend par \guillemetleft \gls{IA} \guillemetright \space la technologie capable d'effectuer des tâches assimilées à l'intelligence comme l'apprentissage et le raisonnement. Cela se concrétise notamment avec le \gls{ML} et le \gls{DPL} permettant une autonomisation de ces outils demandant moins de supervision humaine. L'\gls{IA} est de plus en plus intégrée au monde culturel et patrimonial : d'abord avec la \gls{CV} pour l'analyse d'image et de vidéo et le \gls{NLP} pour le traitement des textes. Puis ensuite avec l'arrivée des \gls{VLM} permettant de faire aussi bien de la \gls{CV} que du \gls{NLP} à l'aide d'un seul outil.
		
		Toutes ces nouvelles technologies se mettent peu à peu au profit du patrimoine a l'instar du projet HIMANIS a permis l'indexation massive des 75 000 pages manuscrites du corpus médiéval du Trésor des Chartes grâce au \gls{HTR}, ou reconnaissance d'écriture manuscrite.\footnote{ Pour plus d'informations voir : \cite{stutzmannRecherchePleinTexte2017/dec./15}}  La conduite ponctuelle de projet d'\gls{IA} apporte plus de profondeur dans notre connaissance du patrimoine possédé. Mais aujourd'hui l'attention se concentre sur l'intégration pérenne d'outils d'\gls{IA}. Des institutions nationales intégrent progressivement de tels outils constituant de véritables processus automatiques : on peut penser à LECTAUREP pour les archives notariales aux Archives Nationales et Gallica Images à la \gls{BNF} pour l'indexation automatique de documents numérisés.  Ces outils sont cependant marginaux puisque les institutions plus petites, mais plus nombreuses, se dirigent vers des solutions de \terme{SGC} sur étagère. Ces dernières intègrent rarement ces outils et seulement de manière partielle avec une prise en main limitée. 
		
		Mon stage porte précisément sur ce sujet avec pour angle particulier de penser les outils et la manière de les intégrer dans une solution de \terme{SGC} développée par l'entreprise SkinSoft à destination de professionnels du monde patrimonial.
		
		Établie à Besançon, SkinSoft est définie comme une PME française, composée d'environ 50 personnes, ainsi qu'un laboratoire de recherche qui fournit une solution de \terme{SGC}. Cette entreprise est une actrice de premier plan sur le marché international des \terme{SGC}, fournissant aussi bien des institutions publiques \textit{(Château de Fontainebleau, The Royal BC Museum, ...)} que privées \textit{(Fondation Jérôme Seydoux-Pathé, Lacoste, ...)}. Elle opère en France principalement, mais aussi en Amérique, en Europe centrale et de l'Est allant même jusqu'en Asie. Au titre de laboratoire de recherche, elle cherche à développer de nouveaux outils pour la gestion des collections notamment basés sur l'\gls{IA}. 
		
		Plusieurs missions ont été définies, chacune portant sur un point d'attention particulier quand on parle d'intégration d'outils \gls{IA} dans un \terme{SGC}. La première a été de réfléchir à ce que l'\gls{IA} peut apporter dans l'application à un utilisateur\footnote{Ici le mot \guillemetleft utilisateur \guillemetright désigne tout professionnel membre d'une institution, privée ou patrimoniale, possédant la solution SkinSoft et qui utilise cette dernière pour gérer la collection de ladite institution}, sur des aspects d'usages et d'accessibilités des collections. La seconde mission a consisté en l'identification des besoins utilisateurs et la conception d'outils y répondant. La troisième et dernière mission a été de conceptualiser une chaîne de traitement basée sur l'\gls{IA}\footnote{Ici, le terme \guillemetleft intelligence artificielle \guillemetright est utilisé au sens global. Nous ne parlons pas encore de domaine, sous-domaine et d'outils précis.} afin d'automatiser le traitement de l'\gls{imgani} directement dans l'application. 
		
		La réalisation de ces missions a pour objectif de produire différents outils basés sur l'\gls{IA} afin de permettre une meilleure gestion des grandes quantités de données patrimoniales. D'offrir de meilleurs outils de gestion et d'exploitation des collections et données. Et enfin de faciliter les processus assurés par un \terme{SGC}, comme l'indexation, notamment dans le champ du traitement de l'\gls{imgani}.\footnote{Il est important de noter que ce travail ne porte pas uniquement sur l'\gls{imgani}. Le travail effectué explore d'autres typologies que la seule \gls{imgani}, bien que cette dernière possède une place particulière ici.}
		\newline
		
		Finalement, cet objectif pose la question suivante : dans quelle mesure l’implémentation d'outils basés sur l’intelligence artificielle dans un CMS permet de répondre aux défis des collections composites sans dénaturer les pratiques de gestion ?\newline
		
		
		%Pourquoi implémenter des outils d’intelligence artificielle dans des systèmes de gestion de collection et comment les concevoir face aux défis de la gestion des collections ainsi que des données patrimoniales publiques et privées ?ple objectif  
		
		La réponse se fait en 3 temps. Le premier est consacré à l'analyse des \terme{SGC}, notamment le système SkinSoft, et des difficultés actuelles qu'ils rencontrent. C'est également l'occasion de quelques considérations générales sur l'\gls{IA} dans le domaine de la gestion des collections. L'objectif est de solidifier nos connaissances sur le monde dans lequel cette étude s'inscrit et de comprendre en profondeur chaque élément d'attention ciblés quand nous parlons de gestion de collections d'\gls{IA} appliquée à ce domaine particulier. 
		
		Le second temps est dédié à l'exploration de deux technologies d'\gls{IA} : en particulier le \gls{NLP} et l'indexation semi-automatisée\footnote{Ce n'est pas vraiment une technologie à part entière, mais plutôt un ensemble de technologies mobilisées pour une tâche particulière. Par facilité de langage, nous parlerons souvent de l'indexation semi-automatisée comme d'une technologie à part entière.}. L'analyse de l'indexation automatisée notamment via une réflexion à son application au domaine de l'\gls{imgani}. Nous ferons un état de l'art afin de comprendre l'envergure de ces technologies, les problématiques qu'elles résolvent et comment la recherche actuelle conçoit leurs utilisations sur des données patrimoniales. La solution SkinSoft sert ici de cadre d'étude permettant de mieux saisir le périmètre et les limites de chaque technologie. 
		
		Le troisième et dernier temps de ce mémoire est dédié à un retour d'expérience sur les deux projets d'\gls{IA}. En présentant les solutions conceptualisées et comment leurs implémentations ont été pensées. L'objectif est de présenter les apports et limites de ces solutions ainsi que les difficultés rencontrées et les choix effectués. Avant de conclure notre travail, ce temps se termine sur une analyse d'écart entre l'apport envisagé par l'implémentation de ces technologies au début des projets, et l'apport réellement envisageable une fois les projets aboutis. 